% Wymaga: booktabs. Nie wymaga pakietu float.
\makeatletter
\@ifundefined{tablehere}{\newenvironment{tablehere}{\par\vskip\intextsep\def\@captype{table}\begin{center}}{\end{center}\par\vskip\intextsep}}{}
\makeatother


\subsection{Optymalizacja sieci i serwera}
\subsubsection{Opis optymalizacji}

W pierwotnej wersji serwisu strona pobierała najpierw listę wpisów (\texttt{GET /api/posts},
zwracającą tylko identyfikatory i tytuły), a następnie pięć kolejnych, wykonywanych
jeden po drugim żądań o szczegóły każdego wpisu (problem N+1). Przy symulowanym opóźnieniu
50--150\,ms na żądanie dawało to ponad 600\,ms łańcucha żądań. Serwer nie kompresował
odpowiedzi, nie wysyłał nagłówków \texttt{Cache-Control} (powtórne odwiedziny
ponownie walidowały lub pobierały zasoby), a przy każdym żądaniu i zdarzeniu
analitycznym wykonywał synchroniczny zapis na dysk (\texttt{appendFileSync}),
który blokował pętlę zdarzeń. Klient wysyłał zdarzenia synchronicznym
\texttt{XMLHttpRequest}, blokującym wątek główny przeglądarki.

Zmiany:
\begin{itemize}
  \item \texttt{GET /api/posts} zwraca pełne wpisy; klient wykonuje jedno żądanie
    zamiast sześciu,
  \item kompresja odpowiedzi tekstowych (Brotli lub gzip, zależnie od
    \texttt{Accept-Encoding}) z użyciem wbudowanego modułu \texttt{zlib},
  \item nagłówki cache'ujące (patrz tabela poniżej),
  \item logi zapisywane przez strumienie \texttt{fs.createWriteStream} w trybie dopisywania,
  \item zdarzenia wysyłane przez \texttt{navigator.sendBeacon} (z \texttt{fetch}
    z \texttt{keepalive} jako zapasowym rozwiązaniem) zamiast synchronicznego XHR,
  \item obraz Dockera: \texttt{node:20-alpine}, instalacja zależności
    (\texttt{npm ci --omit=dev}) przed skopiowaniem kodu oraz \texttt{HEALTHCHECK}.
\end{itemize}

\subsubsection{Konfiguracja i parametry}

\begin{tablehere}
\centering
\begin{tabular}{l p{9cm}}
\toprule
\textbf{Zasób} & \textbf{Nagłówek} \\
\midrule
HTML, własne CSS i JS & \texttt{Cache-Control: no-cache} -- przeglądarka zawsze pyta
  o ważność (ETag), a przy braku zmian dostaje \texttt{304} bez treści; nazwy plików
  nie zawierają skrótu, więc dłuższe cache'owanie mogłoby serwować nieaktualną wersję \\
Biblioteki, obrazy, wideo, czcionki & \texttt{public, max-age=31536000, immutable}
  -- bez ponownych żądań przez rok (zmianę pliku trzeba wykonać pod nową nazwą) \\
\texttt{/api/posts} & \texttt{public, max-age=60, stale-while-revalidate=600} --
  przez minutę odpowiedź z cache, potem do 10 minut serwowana z cache
  przy odświeżaniu w tle \\
\bottomrule
\end{tabular}
\caption{Polityka cache'owania}
\end{tablehere}

\begin{tablehere}
\centering
\begin{tabular}{l p{9cm}}
\toprule
\textbf{Parametr} & \textbf{Znaczenie} \\
\midrule
\texttt{Accept-Encoding} & wybór Brotli (\texttt{br}) lub gzip; bez obsługi -- odpowiedź bez kompresji \\
\texttt{BROTLI\_PARAM\_QUALITY: 4} & niski poziom kompresji Brotli -- dobry kompromis
  między stopniem kompresji a obciążeniem CPU przy kompresji w locie \\
\texttt{Vary: Accept-Encoding} & cache pośredni rozróżnia wersje skompresowane i nie \\
\texttt{COMPRESSIBLE} & kompresowane są tylko typy tekstowe (HTML, CSS, JS, JSON, SVG); obrazy i wideo
  są już skompresowane \\
\texttt{flags: 'a'} & strumień dopisujący do pliku; zapis asynchroniczny \\
\bottomrule
\end{tabular}
\caption{Parametry kompresji i logowania}
\end{tablehere}

Kompresję zaimplementowano własnym oprogramowaniem pośredniczącym na \texttt{zlib},
aby nie dodawać zależności do \texttt{package.json}.

\subsubsection{Wyniki pomiarów}

Pomiary teoretyczne wynikają z liczby żądań i symulowanego opóźnienia; rozmiary skompresowane
nie zostały jeszcze zmierzone (patrz tabela wyników).

\begin{tablehere}
\centering
\begin{tabular}{l r r}
\toprule
\textbf{Wielkość} & \textbf{Przed} & \textbf{Po} \\
\midrule
Żądania API przy ładowaniu & 6 (szeregowo) & 1 \\
Opóźnienie symulowane (średnio) & ok.\ 6\,$\times$\,100\,ms = 600\,ms & ok.\ 100\,ms \\
Zapisy synchroniczne na dysk & 1 na żądanie i zdarzenie & 0 \\
Synchroniczny XHR w przeglądarce & 1 & 0 \\
Kompresja tekstu & brak & br / gzip \\
Powtórna wizyta & ponowny transfer lub walidacja wszystkich zasobów & \texttt{304} dla HTML/CSS/JS, reszta z cache \\
\bottomrule
\end{tabular}
\caption{Zmiany w liczbie żądań i kosztach serwera}
\end{tablehere}

\begin{tablehere}
\centering
\begin{tabular}{l c c c c c}
\toprule
\textbf{Metryka} & \textbf{PSI mobile} & \textbf{PSI desktop} & \textbf{GTmetrix} & \textbf{Zmiana vs. stan wyjściowy} & \textbf{Uwagi} \\
\midrule
LCP [s]        & [—] & [—] & [—] & [—] & [—] \\
TBT [ms]       & [—] & [—] & [—] & [—] & [—] \\
CLS            & [—] & [—] & [—] & [—] & [—] \\
FCP [s]        & [—] & [—] & [—] & [—] & [—] \\
Speed Index [s]& [—] & [—] & [—] & [—] & [—] \\
\midrule
Wynik wydajności (0–100) & [—] & [—] & [—] & [—] & [—] \\
Rozmiar strony [MB]      & —   & —   & [—] & [—] & [—] \\
Liczba żądań HTTP        & —   & —   & [—] & [—] & [—] \\
Transfer przy powtórnej wizycie [KB] & — & — & [—] & [—] & [—] \\
\texttt{/api/posts} p95 [ms] & — & — & [—] & [—] & [—] \\
\bottomrule
\end{tabular}
\caption{Wyniki pomiarów dla optymalizacji sieci i serwera}
\end{tablehere}

\subsubsection{Wnioski}

Największy zysk w czasie ładowania daje usunięcie wzorca N+1: sześć szeregowych
żądań zastąpiono jednym. Kompresja zmniejsza transfer plików tekstowych, a
cache'owanie sprawia, że powtórne odwiedziny prawie nic nie kosztują. Zamiana
zapisów synchronicznych na strumienie i synchronicznego XHR na \texttt{sendBeacon}
usuwa blokowanie pętli zdarzeń serwera oraz wątku głównego przeglądarki.

Ograniczeniem jest brak skrótów w nazwach plików własnego CSS i JS (dlatego
używają walidacji, a nie \texttt{immutable}) oraz to, że zmiana obrazu lub filmu
wymaga zmiany jego nazwy. Dane o rozmiarach po kompresji i czasach należy uzupełnić
po pomiarach. Zmiany nie obejmują samych mediów ani bibliotek front-endowych --
to osobne optymalizacje.
